\documentclass[11pt]{article}
\usepackage[a4paper,margin=2.5cm]{geometry}
\usepackage{amsmath,amssymb}
\usepackage{booktabs}
\usepackage{longtable}
\usepackage{tabularx}
\usepackage{array}
\usepackage{siunitx}
\usepackage[hidelinks]{hyperref}

\title{Appendix: Extraction of the energetic value of composite diets}
\author{}
\date{}

\begin{document}
\maketitle

\section*{Theory}

\subsection*{DEB Formulation for a single compound X in a substrate}

In DEB theory, compounds are quantified on a carbon-molar basis to ensure consistency between chemical composition, thermodynamic potential, and energetic bookkeeping. For a compound $X$ with empirical formula $\mathrm{C}_{n_C}\mathrm{H}_{n_H}\mathrm{O}_{n_O}\mathrm{N}_{n_N}$, the carbon content defines the fundamental counting unit: one C-mol (carbon mole) is the amount of compound that contains one mole of carbon atoms. Thus, $1\,\mathrm{mol}$ of compound contains $n_C$ C-moles. All concentrations refer to the substrate in which compound $X$ is present.

Following DEB convention, the carbon-molar concentration is denoted by $X$, while the corresponding mass concentration of compound $X$ in the substrate is denoted by $c_X$ ($\mathrm{g\,cm^{-3}}$). The relations between chemical composition, concentration, chemical potential, and resulting energy density due exclusively to compound $X$ are summarized in Table~\ref{tab:single-compound}.

\begin{table}[h]
\centering
\small
\begin{tabularx}{\textwidth}{>{\raggedright\arraybackslash}p{1.8cm} >{\raggedright\arraybackslash}p{4.4cm} >{\raggedright\arraybackslash}p{2.7cm} >{\raggedright\arraybackslash}X}
\toprule
Symbol & Description & Unit & Defining equation \\
\midrule
$w$ & Molar mass & $\mathrm{g\,mol^{-1}}$ & $w = 12 n_C + 1 n_H + 16 n_O + 14 n_N$ \\
$w_C$ & Molar mass per C-mol & $\mathrm{g\,C\text{-}mol^{-1}}$ & $w_C = \dfrac{w}{n_C}$ \\
$c_X$ & Mass concentration in the substrate & $\mathrm{g\,cm^{-3}}$ & --- \\
$X$ & Carbon-molar concentration in the substrate & $\mathrm{C\text{-}mol\,cm^{-3}}$ & $X = \dfrac{c_X}{w_C}$ \\
$\mu_X$ & Chemical potential (per C-mol) & $\mathrm{J\,C\text{-}mol^{-1}}$ & --- \\
$e_X$ & Energy per gram & $\mathrm{J\,g^{-1}}$ & $e_X = \dfrac{\mu_X}{w_C}$ \\
$\varepsilon_X$ & Energy density in the substrate & $\mathrm{J\,cm^{-3}}$ & $\varepsilon_X = \mu_X X = e_X c_X$ \\
\bottomrule
\end{tabularx}
\normalsize
\caption{Definitions and relations for compound $X$ in a substrate under the carbon-molar formulation of DEB theory. Energy densities refer exclusively to the energetic contribution of compound $X$.}
\label{tab:single-compound}
\end{table}

\subsection*{Major Polymer Classes in DEB Theory}

In DEB theory, structural biomass and reserves are represented as mixtures of three principal polymer groups: carbohydrates (polysaccharides), lipids, and proteins. Their average empirical compositions and energetic properties are given in Table 4.2 of the DEB textbook (combustion reference, pH 7, 25 $^\circ$C). The empirical formula is expressed on a per C-mol basis, i.e. $n_C=1$.

\begin{table}[h]
\centering
\small
\begin{tabularx}{\textwidth}{>{\raggedright\arraybackslash}p{3.0cm} >{\raggedright\arraybackslash}p{3.2cm} >{\raggedright\arraybackslash}p{2.8cm} >{\raggedright\arraybackslash}p{2.8cm} >{\raggedright\arraybackslash}X}
\toprule
Compound & Empirical formula $\mathrm{C}\mathrm{H}_{n_H}\mathrm{O}_{n_O}\mathrm{N}_{n_N}$ & Molar mass per C-mol $w_{C,X}$ ($\mathrm{g\,C\text{-}mol^{-1}}$) & Chemical potential $\mu_X$ ($\mathrm{kJ\,C\text{-}mol^{-1}}$) & Energy per gram $e_X$ ($\mathrm{kJ\,g^{-1}}$) \\
\midrule
Carbohydrates (polysaccharides) & $\mathrm{CH_2O}$ & 30.00 & 516 & 17.2 \\
Lipids & $\mathrm{CH_{1.92}O_{0.12}}$ & 15.84 & 616 & 38.9 \\
Proteins & $\mathrm{CH_{1.61}O_{0.33}N_{0.28}}$ & 22.81 & 401 & 17.6 \\
\bottomrule
\end{tabularx}
\normalsize
\caption{Average empirical composition and energetic properties of the three principal polymer classes in DEB theory (combustion reference frame).}
\label{tab:polymer-classes}
\end{table}

\subsection*{Frequently used dietary compounds in \textit{Drosophila} media}

Edible components of \textit{Drosophila} diet media are characterized by their empirical formula and their chemical potential $\mu_X$, from which their C-molar mass $w_{C,X}$ and energetic variables $e_X,\varepsilon_X$ can be derived (see Table~\ref{tab:single-compound}). Non-edible structural components of the substrate (e.g., agar) are excluded since they do not contribute metabolizable energy. The most commonly used energetically meaningful compounds are shown in Table~\ref{tab:drosophila-compounds}.

\begin{table}[h]
\centering
\small
\begin{tabularx}{\textwidth}{>{\raggedright\arraybackslash}p{3.0cm} >{\raggedright\arraybackslash}p{3.6cm} >{\raggedright\arraybackslash}p{2.0cm} >{\raggedright\arraybackslash}p{2.6cm} >{\raggedright\arraybackslash}X}
\toprule
Compound & Empirical formula $\mathrm{C}_{n_C}\mathrm{H}_{n_H}\mathrm{O}_{n_O}\mathrm{N}_{n_N}$ & Molar mass per C-mol $w_{C,X}$ ($\mathrm{g\,C\text{-}mol^{-1}}$) & Chemical potential $\mu_X$ ($\mathrm{kJ\,C\text{-}mol^{-1}}$) & Energy per gram $e_X$ ($\mathrm{kJ\,g^{-1}}$) \\
\midrule
Glucose / Fructose / Dextrose & $\mathrm{C}_6\mathrm{H}_{12}\mathrm{O}_6$ & 30.03 & 468 & 15.6 \\
Sucrose & $\mathrm{C}_{12}\mathrm{H}_{22}\mathrm{O}_{11}$ & 28.53 & 471 & 16.5 \\
Starch (repeat unit) & $(\mathrm{C}_6\mathrm{H}_{10}\mathrm{O}_5)_n$ & 27.02 & 459 & 17.0 \\
Brewer's yeast (dry) & $\approx \mathrm{C}_{4.5}\mathrm{H}_{7.3}\mathrm{O}_{2.2}\mathrm{N}_{0.8}$ & $\approx 25.11$ & 478 & 19.0 \\
Propionic acid & $\mathrm{C}_3\mathrm{H}_6\mathrm{O}_2$ & 24.69 & 513 & 20.8 \\
\bottomrule
\end{tabularx}
\normalsize
\caption{Chemical composition and energetic properties of edible compounds used in \textit{Drosophila} diet media.}
\label{tab:drosophila-compounds}
\end{table}

\subsection*{Computation of energy density of a composite diet}

In practice, for a batch of diet $X$ of volume $V$ (cm$^3$) containing assimilable ingredients $i$, recipe information can be written in different but equivalent forms:
\begin{itemize}
  \item If ingredient masses $m_i$ (g) are known, the mass concentration of ingredient $i$ is
  \begin{equation}
    c_i = \frac{m_i}{V}\quad [\mathrm{g\,cm^{-3}}].
  \end{equation}
  \item If ingredient $i$ is given as a mixture fraction (e.g., yeast paste), with mass fraction $w_i$ in an aqueous/jelly medium, then
  \begin{equation*}
    c_i = w_i\,\rho_{\text{medium}}.
  \end{equation*}
  We assume for slurry-like aqueous or jelly media that $\rho_{\text{medium}}\approx 1\,\mathrm{g\,cm^{-3}}$.
  \item If ingredient $i$ is given as $w/V$ (e.g., glucose--agarose), then $c_i$ follows directly from Eq.~(1), i.e.
  \begin{equation*}
    c_i = \frac{m_i}{V}
    \quad\text{(for example, }10\%\,w/v = 10\,\mathrm{g}/100\,\mathrm{cm^3} = 0.10\,\mathrm{g\,cm^{-3}}\text{).}
  \end{equation*}
\end{itemize}



Once we have the recipe of the diet, with known values of $c_i$, $e_i$, and $w_{C,i}$ for all ingredients (with $e_i=\mu_i/w_{C,i}$), we can compute the composite energy density, DEB food density, and effective chemical potential as:
\begin{equation}
  \varepsilon_X = \sum_i c_i e_i\quad [\mathrm{J\,cm^{-3}}].
\end{equation}
\begin{equation}
  X = \sum_i \frac{c_i}{w_{C,i}}\quad [\mathrm{C\text{-}mol\,cm^{-3}}].
\end{equation}
\begin{equation}
  \mu_X = \frac{\varepsilon_X}{X}\quad [\mathrm{J\,C\text{-}mol^{-1}}].
\end{equation}


\section*{Example diets from the literature}

Below we apply the above theory to extract DEB food descriptors for the diets used in the literature. 
For every study, for each diet, we first compute the recipe in mass concentrations $c_i$ (g~cm$^{-3}$) of the assimilable ingredients.
The molar mass and energy per gram for each ingedient are derived from the compound data in Table~\ref{tab:polymer-classes} and Table~\ref{tab:drosophila-compounds}. 
Then we compute the energy density $\varepsilon_X$, food density $X$, and chemical potential $\mu_X$ for the composite diet using Eq.~(1), Eq.~(2), and Eq.~(3). 

\section*{Study I : Kaun et al. (2007)}

In Kaun et al.\ (2007), larvae were reared on a standard culture medium 
containing yeast and sucrose as the principal nutritive components 
(together with agar and mineral salts), and food-quality treatments were 
generated by proportionally reducing both yeast and sucrose to 75\%, 50\%, 
25\%, or 15\% of this formulation (Materials and methods, p.~2). This medium is therefore a multi-compound nutritive substrate. 

For intake assays, larvae were transferred to yeast paste (2:1 water:yeast) containing dye or radiolabeled substrates. In the $^{14}$C-glucose assays, labelled glucose is treated only as a tracer and is excluded from DEB energy accounting; thus, in energetic terms, yeast paste is a single-compound diet (yeast biomass). The yeast paste is prepared as 2:1 water:yeast (w/w), so the yeast mass fraction is $w_Y=1/3$. Assuming a slurry density $\rho_{\text{paste}}\approx 1\,\mathrm{g\,cm^{-3}}$, the yeast mass concentration is $c_Y = w_Y\,\rho_{\text{paste}} = 0.333\,\mathrm{g\,cm^{-3}}$.

The glucose--agarose and fructose--agarose substrates are also single-compound diets in energetic terms, because each contains one nutritive compound dissolved in agarose. For the sugar--agarose diets, the glucose--agarose recipe (10\% glucose w/v, 2.3\% agarose) gives $c_G=0.10\,\mathrm{g\,cm^{-3}}$, while the fructose--agarose recipe (2 mol\,L$^{-1}$ fructose, 1\% agarose) gives
\begin{equation*}
  c_F = \frac{2\times 180.16}{1000} = 0.36032\,\mathrm{g\,cm^{-3}}.
\end{equation*}

Only the standard rearing medium is multi-compound in this study.



The assimilable recipe concentrations for all diets are summarized in Table~\ref{tab:study1-recipes}.

\begin{table}[h]
\centering
\small
\begin{tabularx}{\textwidth}{>{\raggedright\arraybackslash}p{3.2cm} >{\centering\arraybackslash}p{2.2cm} >{\centering\arraybackslash}p{2.2cm} >{\centering\arraybackslash}p{2.2cm} >{\centering\arraybackslash}X}
\toprule
Diet & Yeast & Sucrose & Glucose & Fructose \\
\midrule
Standard (100\%) & 0.0500 & 0.1000 & 0 & 0 \\
Standard (quality $q$) & $q\times 0.0500$ & $q\times 0.1000$ & 0 & 0 \\
Yeast paste & 0.3330 & 0 & 0 & 0 \\
Glucose--agarose & 0 & 0 & 0.1000 & 0 \\
Fructose--agarose & 0 & 0 & 0 & 0.3603 \\
\bottomrule
\end{tabularx}
\normalsize
\caption{Assimilable-compound recipes for diets in Kaun et al. (2007), expressed as mass concentrations $c_i$ in g~cm$^{-3}$.}
\label{tab:study1-recipes}
\end{table}

We now apply the formulas to these diet recipes and compute the energetic variables for each diet; for diluted food with quality factor $q$ (e.g. 0.75, 0.50, 0.25, 0.15), which scales yeast and sucrose together, $\varepsilon(q)=q\,\varepsilon_{\text{standard}}$ and $X(q)=q\,X_{\text{standard}}$. The results are summarized in Table~\ref{tab:summary}.


\begin{table}[h]
\centering
\small
\begin{tabularx}{\textwidth}{>{\raggedright\arraybackslash}p{2.8cm} >{\raggedright\arraybackslash}p{2.8cm} >{\raggedright\arraybackslash}p{2.8cm} >{\raggedright\arraybackslash}p{3.0cm} >{\raggedright\arraybackslash}X}
\toprule
Diet & Energy density, $\varepsilon_X$ (J\,cm$^{-3}$) & Food density, $X$ (C-mol\,cm$^{-3}$) & Chemical potential, $\mu_X$ (J\,C-mol$^{-1}$) & Notes \\
\midrule
Standard (100\%) & $2.7\times10^{3}$ & $5.50\times10^{-3}$ & $4.91\times10^{5}$ & complete diet \\
Standard (quality $q$) & $q\times 2.7\times10^{3}$ & $q\times 5.50\times10^{-3}$ & $4.91\times10^{5}$ & dilution series \\
Yeast paste & $7.0\times10^{3}$ & $1.33\times10^{-2}$ & $5.28\times10^{5}$ & high-quality \\
Glucose--agarose & $1.56\times10^{3}$ & $3.33\times10^{-3}$ & $4.68\times10^{5}$ & carbon-only test \\
Fructose--agarose & $5.62\times10^{3}$ & $1.20\times10^{-2}$ & $4.68\times10^{5}$ & carbon-only test \\
\bottomrule
\end{tabularx}
\normalsize
\caption{Energetic variables for diets in Kaun et al. (2007), computed from the assimilable-compound recipe concentrations in Table~\ref{tab:study1-recipes}.}
\label{tab:summary}
\end{table}

\section*{Study II: Ormerod et al. (2017)}

From Table~2 of Ormerod et al. (2017), wet percentages are treated as g~cm$^{-3}$ assuming bulk density $\approx 1$~g~cm$^{-3}$.

\begin{table}[h]
  \centering
  \sisetup{table-number-alignment = center}
  \begin{tabular}{lS[table-format=1.4]S[table-format=1.4]S[table-format=1.4]}
    \toprule
    Diet & {Protein} & {Fat} & {Carbohydrate} \\
    \midrule
    F424 & 0.0191 & 0.0012 & 0.2119 \\
    JAZZ & 0.0080 & 0.0001 & 0.1701 \\
    \bottomrule
  \end{tabular}
  \caption{Assimilable-compound recipes for diets in Ormerod et al. (2017), expressed as mass concentrations $c_i$ in g~cm$^{-3}$.}
\end{table}

results

\begin{table}[h]
  \centering
  \sisetup{table-number-alignment = center}
  \begin{tabular}{lS[table-format=1.2e1]S[table-format=1.2e-2]S[table-format=1.2e5]}
    \toprule
    Diet & {$\varepsilon_X$ (J~cm$^{-3}$)} & {$X$ (C-mol~cm$^{-3}$)} & {$\mu_X$ (J~C-mol$^{-1}$)} \\
    \midrule
    \textbf{F424} & 4.03e3 & 7.98e-3 & 5.05e5 \\
    \textbf{JAZZ} & 3.07e3 & 6.03e-3 & 5.09e5 \\
    \bottomrule
  \end{tabular}
  \caption{Absolute DEB food descriptors for the two diets.}
\end{table}

\section*{Study III: Ng’oma et al. (2019)}

\end{document}